\begin{prob}
    What is the \emph{best} case time complexity of \python{quickselect}?
    Describe why and provide the input that yields this best case.
    For simplicity, you may assume that Quickselect always chooses the last element of the input
    array as the pivot (that is, the pivot is \emph{not} chosen randomly).

    \begin{soln}
        The best case time complexity of quickselect is $\Theta(n)$, where $n$ is the
        size of the input array.

        In the best case, the element chosen as the first pivot is exactly the
        order statistic being queried so that no recursive calls are made.  However, we still must perform the
        partition operation to count the number of elements to the left and
        right of the pivot, and this takes linear time.

        Alternatively, another case which also leads to $\Theta(n)$ run time (and so is
        also a ``best case'') is when the chosen pivot is not the order statistic we are
        looking for, but it is the median (or close to it). Then there is still a recursive
        call on an array of size $n/2$, plus $\Theta(n)$ work for the partition. If the
        recursive calls also choose the median as their pivot, this gives a recurrence of
        $T(n) = T(n/2) + n$. Solving this gives a time complexity of $\Theta(n)$.
    \end{soln}
\end{prob}
